\begin{textsample}{POS Dim 3 – human – Score 6.00 – rj/rj\_inf\_e\_ef\_intro\_4\_p3.txt}  \label{ex:f3_pos_012}
Essas decisões precisam, igualmente, ser consideradas na organização de currículos e \textbf{propostas} adequados às diferentes modalidades de ensino ( Educação Especial, Educação de Jovens e Adultos, Educação do Campo, Educação Escolar Indígena, Educação Escolar Quilombola, Educação a Distância ), atendendo -se às orientações das \textbf{Diretrizes} \textbf{Curriculares} Nacionais.
Faz -se necessário então, tendo este \textbf{Documento} de Orientação \textbf{Curricular} como referência, que as redes de ensino formulem \textbf{propostas} \textbf{curriculares} para as modalidades de ensino acima citadas considerando as \textbf{especificidades} requeridas por estes estudantes de maneira a prover condições adequadas de progressão de aprendizagem para estes públicos.

% matched lemmas: curricular, diretriz, documento, especificidade, proposta
\end{textsample}
